% =====================================================================
%  General introduction
% =====================================================================
\chapter*{General Introduction}
\addcontentsline{toc}{chapter}{General Introduction}
\markboth{GENERAL INTRODUCTION}{GENERAL INTRODUCTION}

The growing affordability of GPS tracking devices, the maturity of cloud platforms and the widespread adoption of telematics solutions across logistics fleets have produced an environment in which raw vehicular data is no longer the bottleneck. The main challenge is no longer data collection. It is the design of analytical pipelines able to transform large volumes of telemetry into reliable and actionable decision-support indicators. Within this context, fleet operators face two parallel pressures: the need to reduce operational risks (accidents, fuel waste, maintenance costs) and the need to provide their operations teams with explanatory dashboards that translate technical metrics into management decisions. Beyond the descriptive layer that the industry has now mastered, value increasingly comes from \emph{advanced analytics} --- the disciplined transformation of raw and archived data into segmentations, anomaly scores and risk indicators that support faster, evidence-based decisions, allow abnormal patterns to be identified early and used as risk-prioritization signals, and feed strategic planning rather than only post-hoc reporting.

The host organization of this thesis operates a multi-tenant telematics platform that supports a sizeable fleet of instrumented vehicles across several customer accounts. Over years of operation, the company has accumulated a substantial historical archive of decoded telemetry alongside its live operational stack. Its current analytical surface --- composed of dashboards and standard reports --- meets the descriptive needs of fleet managers but stops short of forward-looking signals. What is missing is not a visualization layer; it is the advanced analytical layer that would turn the historical archive into actionable behavioural and risk indicators. The internship that this thesis documents was scoped explicitly to fill that gap.

The objective of the project is to design and to industrialize a CRISP-DM-aligned analytical platform that migrates the relevant slices of the operational data into a dedicated analytical layer, materializes a clean warehouse and a family of analytical marts, trains an unsupervised behavioural segmentation at the device-month grain, evaluates an anomaly-oriented risk indicator and exposes a transparent risk profile through an Application Programming Interface (API) and a Business Intelligence (BI) dashboard. The technical motivation for the chosen analytical Database Management System (DBMS), the detailed presentation of the host organization and its operational data sources, and the collaborative working setup adopted for the internship are all developed in the dedicated chapters of the report.

\paragraph{Main contributions.} The work documented in this thesis delivers five concrete contributions that go beyond the descriptive analytical surface of the existing platform:
\begin{enumerate}
    \item a \textbf{MySQL-to-PostgreSQL conversion bridge} that ingests the operational dump slices and stages them in a typed analytical schema;
    \item an \textbf{analytical PostgreSQL warehouse} consolidating decoded telemetry and master data into a query-friendly model dedicated to advanced analytics;
    \item a \textbf{transparent cleaning rule engine} and a \textbf{device-month mart} that turn cleaned telemetry into behavioural and exposure indicators suitable for modeling;
    \item an \textbf{unsupervised behavioural segmentation} of the fleet and an \textbf{anomaly-driven risk score} at the device-month grain, accompanied by per-tenant clustering quality controls;
    \item the \textbf{deployment of an API and a BI dashboard} that expose the risk profile to operational users, supported by cron-based automation and a basic monitoring layer.
\end{enumerate}

The thesis is organized in six chapters. \textbf{Chapter~\ref{chap:context}} situates the project in its industrial and academic context, presents the host organization and its three operational MySQL databases, analyses the existing analytical surface and motivates the choice of CRISP-DM as the guiding methodology. \textbf{Chapter~\ref{chap:business}} formalizes the business understanding, lists the research questions, derives the functional and non-functional requirements, and presents the target architecture and technology stack --- including the technical justification of the analytical DBMS and the conversion bridge towards it. \textbf{Chapter~\ref{chap:data}} documents the data understanding phase, with schema inspection, data dictionary, missing-value analysis and exploratory data analysis on the staging layer derived from the converted dumps. \textbf{Chapter~\ref{chap:prep}} describes the conversion bridge, the data preparation pipeline, the cleaning rule engine and the engineering of the analytical features. \textbf{Chapter~\ref{chap:model}} combines the CRISP-DM modeling and evaluation phases, explains the reformulation from supervised driver scoring to device-month unsupervised modeling, compares candidate techniques and distinguishes research artefacts from deployed scoring contracts. \textbf{Chapter~\ref{chap:deploy}} presents the deployment of the pipeline and the models in the shared analytical environment, the API and the dashboard, the cron-based automation, the monitoring layer and the secure remote operations setup. The general conclusion synthesizes the deliverables, the lessons learned and the perspectives opened by the project.

\cleardoublepage
